\section*{Tài liệu tham khảo}
\begin{enumerate}
    \item \url{https://www.zipitwireless.com/blog/4-layers-of-iot-architecture-explained}
    \item \url{https://www.zipitwireless.com/blog/what-is-an-iot-gateway-and-how-does-it-work}
    \item \url{https://www.embitel.com/blog/embedded-blog/understanding-how-an-iot-gateway-architecture-works}
    \item \url{https://m5stack.com/}
\end{enumerate}

\section*{Giao tiếp sử dụng nhiều nhất}
Theo khảo sát sản phẩm, \textbf{Ethernet} và \textbf{Wi-Fi} là các giao tiếp được sử dụng phổ biến nhất trong các gateway IoT thương mại.

\section*{Ví dụ:}
\begin{itemize}
    \item \url{https://www.digikey.com/en/products/detail/lantronix-inc/SGX5150122US/9953933}
    \item \url{https://maplesystems.com/product/modelname/cMT-G02/}
    \item \url{https://docs.m5stack.com/en/core/M5CoreMP135}
\end{itemize}

\section*{Xác định cách cấu hình gateway}
\begin{itemize}
    \item Cần xác định cách cấu hình gateway sao cho khi cắm bất kỳ module nào vào, hệ thống có thể tự nhận diện và áp dụng cấu hình phù hợp.
    \item Các module có thể bao gồm:
    \begin{itemize}
        \item Module Wi-Fi
        \item Module Bluetooth
        \item Module Zigbee
        \item Module LoRa
        \item Module 6LoWPAN
        \item module NB-IoT
    \end{itemize}
    \item Một số nguồn tham khảo:
    \begin{itemize}
        \item \url{https://www.ti.com/lit/wp/spmy013/spmy013.pdf?ts=1752155396994&ref_url=https%253A%252F%252Fwww.google.com%252F}
        \item \url{https://www.packtpub.com/en-us/learning/how-to-tutorials/run-and-configure-iot-gateway}
        \item \url{https://iot-analytics.com/top-10-iot-use-cases/}
        \item \url{https://research.aimultiple.com/iot-applications/}
        \item \url{https://www.particle.io/iot-guides-and-resources/iot-gateway/} - trang này hay phết.
        \item \url{https://www.kaaiot.com/iot-knowledge-base/top-17-protocols-iot-use-cases}
        \item \url{https://www.azilen.com/blog/iot-protocols/}
        \item \url{https://www.pusr.com/blog/Six-types-of-iot-gateway-selection-functions-protocols-and-application-scenarios}
        \item \url{https://thingsboard.io/docs/iot-gateway/}
        \item \url{https://www.slideshare.net/slideshow/iot-hardware-innovation/81363705}
        \item \url{}
        \item \url{}
    \end{itemize}
    \item Cách cấu hình:
    \begin{itemize}
        \item Cấu hình thông qua giao diện người dùng (GUI) hoặc dòng lệnh (CLI). (chưa rõ)
    \end{itemize}
\end{itemize}

\section*{Note phần cần tìm hiểu}
\begin{itemize}
    \item Liên quan đến các từ khóa:
    \begin{itemize}
        \item Linux / MCU - FreeRTOS (Embed OS - ARM)
        \item Low level, Communication protocol
        \item Wire/Wireless Interface to communicate with node/cloud (physical connection)
        \item Communication Protocol / Communication Standard
        \begin{itemize}
            \item Physical layer (frequency, radio wave)
            \item Link layer (frame format)
            \item Protocol requirements
        \end{itemize}
    \end{itemize}
    \item Các giao thức giao tiếp với cloud:
    \begin{itemize}
        \item MQTT
        \item CoAP
        \item HTTP/HTTPS
        \item WebSocket
        \item AMQP
        \item XMPP
    \end{itemize}
    \item Cần tìm hiểu các use case của IoT Gateway.
    \begin{itemize}
        \item các giao thức để giao tiếp giữa các node/end device với gateway:
        \subsection*{1. Giao Thức Truyền Thông Dữ Liệu (Data Link/Protocol Layer)}
        \begin{itemize}
            \item \textbf{Modbus (RTU/ASCII/TCP):}
            \begin{itemize}
                \item Rất phổ biến trong công nghiệp, dùng để truyền dữ liệu giữa gateway và thiết bị qua RS485, RS232 hoặc TCP/IP.
                \item Hỗ trợ cả cảm biến và chấp hành, dễ cấu hình, mạnh về khả năng mở rộng.
            \end{itemize}
            \item \textbf{CAN/CANopen (Controller Area Network):}
            \begin{itemize}
                \item Được dùng nhiều trong ô tô, máy công nghiệp, cho phép nhiều thiết bị giao tiếp song song trên một bus vật lý.
                \item Chuẩn SAE J1939 là biến thể phổ biến trong công nghiệp nặng.
            \end{itemize}
            \item \textbf{IO-Link:}
            \begin{itemize}
                \item Chuẩn mở theo IEC 61131-9, cho phép giao tiếp hai chiều giữa master (gateway) và sensor/actuator.
                \item Hỗ trợ cấu hình, chẩn đoán thiết bị từ xa, thường dùng trong tự động hóa nhà máy.
            \end{itemize}
            \item \textbf{1-Wire, Dallas/Maxim:}
            \begin{itemize}
                \item Giao tiếp đơn giản, dùng cho cảm biến nhiệt độ, độ ẩm, v.v.
            \end{itemize}
        \end{itemize}
        \subsection*{2. Giao Thức Không Dây (Wireless Protocols)}
        \begin{itemize}
            \item \textbf{Zigbee, Z-Wave, Bluetooth Low Energy (BLE):}
            \begin{itemize}
                \item Dùng cho các mạng cảm biến không dây, smart home, smart building.
                \item Gateway cần có module radio và stack phần mềm tương ứng để giao tiếp với thiết bị.
            \end{itemize}
            \item \textbf{LoRa/LoRaWAN:}
            \begin{itemize}
                \item Phù hợp cho các ứng dụng IoT diện rộng, tiết kiệm năng lượng, truyền dữ liệu tầm xa.
            \end{itemize}
            \item \textbf{Wi-Fi:}
            \begin{itemize}
                \item Một số cảm biến/actuator hỗ trợ Wi-Fi, giao tiếp qua TCP/UDP hoặc HTTP/MQTT.
            \end{itemize}
        \end{itemize}
        \subsection*{3. Giao Thức Ứng Dụng (Application Protocols)}
        \begin{itemize}
            \item \textbf{AT Command (tập lệnh AT):}
            \begin{itemize}
                \item Nhiều module modem, cảm biến dùng tập lệnh AT để cấu hình, điều khiển, truy vấn trạng thái.
                \item Gateway phải tích hợp parser để gửi/nhận và xử lý tập lệnh AT.
            \end{itemize}
            \item \textbf{Custom Protocols:}
            \begin{itemize}
                \item Một số thiết bị sử dụng giao thức tự định nghĩa, cần mapping hoặc scripting để gateway hiểu và chuyển đổi dữ liệu.
            \end{itemize}
            \item \textbf{MQTT, CoAP, HTTP (nếu cảm biến thông minh):}
            \begin{itemize}
                \item Một số thiết bị hiện đại tích hợp sẵn giao thức tầng ứng dụng, cho phép giao tiếp trực tiếp với gateway qua mạng IP.
            \end{itemize}
        \end{itemize}
        \begin{itemize}
            \item Kết nối các thiết bị IoT với mạng nội bộ hoặc internet.
        \end{itemize}
        \item các yêu cầu người dùng khác nhau, tùy use case:
        \item các yêu cầu về giao thức để giao tiếp với server/cloud:
        \begin{itemize}
            \item Gateway hỗ trợ giao tiếp với server nào, địa chỉ nào(Ipv4, Ipv6)?
            \item Sử dụng giao thức tầng ứng dụng (HTTP, MQTT...) hay tầng network (UDP, TCP...)?
            \item Có hỗ trợ giao thức bảo mật (SSL/TLS...) không?
        \end{itemize}
        \item các yêu cầu về khả năng cấu hình và mở rộng của gateway:
        \begin{itemize}
            \item Khi cấu hình gateway bằng máy tính:
            \begin{itemize}
                \item Cần có giao diện người dùng (GUI) hay dòng lệnh (CLI)?
                \item Có cần hỗ trợ cấu hình từ xa không?
                \item Có cần hỗ trợ cập nhật firmware qua mạng không?
                \item Cắm vào vị trí nào trên mạch?
                \item Luồng dữ liệu chạy như thế nào trong hệ thống khi cấu hình và khi vận hành? 
                \item Cần xác định cách cấu hình gateway sao cho khi cắm bất kỳ module nào vào, 
                hệ thống có thể tự nhận diện và áp dụng cấu hình phù hợp.
            \end{itemize}
            \item Nếu gặp thiết bị có tập lệnh tự định nghĩa, gateway sẽ hiểu và giao tiếp với
            chúng như thế nào? Có thể định nghĩa tập lệnh mới mà không phải viết lại chương trình không?
            \item Định nghĩa một tập lệnh riêng, buộc thiết bị phải tuân thủ?
            \item Hỗ trợ một số tập lệnh tiêu chuẩn, bỏ qua thiết bị không theo chuẩn?
            \item Kết hợp cả hai hướng trên?
        \end{itemize}
    \item hướng tiếp cận:
        \begin{itemize}
            \item Nhu cầu và vấn đề của người dùng là gì? (phải xác định rõ).
            \item Dự định chọn hướng nào? Lý do?
            \item Định nghĩa sản phẩm dựa trên nhu cầu và vấn đề đó, không chỉ dựa vào kỹ thuật.
            \item Định nghĩa sản phẩm sẽ quyết định thiết kế phần cứng.
            \item Xác định ai là người dùng cuối của sản phẩm.
        \end{itemize}
    \end{itemize}
    \item => tóm lại cần tìm xem thực sự cần một gateway 
    được định nghĩa như thế nào sao cho phù hợp với dự án 
    mình đang làm, kg phải định nghĩa theo kĩ thuật của nó, 
    từ đó mới xác đinh được hướng design hardware cho sản phẩm.
\end{itemize}
